\begin{abstract}
Precise developmental staging is fundamental to biomedical research, yet the pig model---the closest non-primate approximation to human physiology---lacks a molecular definition of development. Current staging relies on subjective morphological criteria that fail to capture the continuous, asynchronous nature of biological maturation. Here, we present a calibrated transcriptomic atlas that systematically maps porcine development across five major tissues (Brain, Liver, Muscle, Lung, and Blood). By integrating large-scale RNA-seq profiles (1,924 samples) with machine learning calibration, our framework achieves high-precision staging (balanced accuracy: 0.64--0.92, mean: 0.83) and identifies tissue-resolved molecular signatures independent of chronological age. Crucially, we provide quantitative evidence that developmental programs in muscle tissue are evolutionarily conserved with humans: cross-species analysis reveals 90\% directional concordance (32/36 genes) and significant correlation (Pearson R = 0.62, p < 10$^{-4}$) in critical neuromuscular and metabolic programs. Key drivers of synaptic maturation (\textit{SORCS2}) and fiber-type specialization (\textit{FGFRL1}) show synchronized trajectories between pig and human muscle. This molecular staging system provides an objective, scalable standard for translational research, enabling precise synchronization of experimental cohorts and confirming the pig's validity as a model for human muscle developmental biology.
\end{abstract}